\chapter{Fourier Analysis}
\label{ch:fourier}

\section{Fourier Series}
\label{sec:fourier-series}

A Fourier series represents any reasonably well--behaved \(2L\)--periodic function \(f(x)\) as an infinite linear combination of sines and cosines.

\subsection{Orthogonal Basis of Sines and Cosines}
The key idea is that the set
\begin{equation}
    \mathcal{S} := \left\{1,\cos\tfrac{\pi x}{L},\sin\tfrac{\pi x}{L},\cos\tfrac{2\pi x}{L},\sin\tfrac{2\pi x}{L},\dots \right\}
\end{equation}
forms an orthogonal basis on the interval \([-L,L]\).

Because sines and cosines of different frequencies are orthogonal, cross--terms disappear when computing the coefficients:

\begin{proposition}{Orthogonality of Sines and Cosines}{orthogonality}
    For integers \( m, n \geq 0 \) on the interval \([-L, L]\), the following orthogonality relations hold:
    \begin{align*}
        \int_{-L}^{L} \cos\left(\tfrac{m\pi x}{L}\right) \cos\left(\tfrac{n\pi x}{L}\right) \, dx
         & =
        \begin{cases}
            0,  & m \neq n,     \\
            L,  & m = n \neq 0, \\
            2L, & m = n = 0,
        \end{cases}             \\
        \int_{-L}^{L} \sin\left(\tfrac{m\pi x}{L}\right) \sin\left(\tfrac{n\pi x}{L}\right) \, dx
         & =
        \begin{cases}
            0, & m \neq n,     \\
            L, & m = n \neq 0,
        \end{cases}              \\
        \int_{-L}^{L} \cos\left(\tfrac{m\pi x}{L}\right) \sin\left(\tfrac{n\pi x}{L}\right) \, dx
         & = 0 \qquad \text{for all } m, n.
    \end{align*}
\end{proposition}

These relations show that the set of functions \( \{1, \cos(\frac{n\pi x}{L}), \sin(\frac{n\pi x}{L}) \} \) forms an orthogonal basis for \( L^2([-L, L]) \).

\subsection{Fourier Series Definition}
If we project \(f\) onto these basis functions, we obtain its Fourier expansion. The coefficients are computed by taking the inner product of \(f\) with each basis function, exploiting their orthogonality.

\begin{definition}{Fourier Series}{fourier-series}
    Let \(f\) be a \(2L\)--periodic function integrable on \([-L,L]\).  Its Fourier series is
    \begin{equation}
        f(x) = \frac{a_0}{2} +\sum_{n=1}^{\infty} \left(a_n\cos\tfrac{n\pi x}{L} +b_n\sin \tfrac{n\pi x}{L}\right)
    \end{equation}

    where the coefficients are the orthogonal projections
    \begin{align*}
        a_0 & =\frac{1}{L}\int_{-L}^{L}f(x)\, dx                      \\
        a_n & =\frac{1}{L}\int_{-L}^{L}f(x)\cos\tfrac{n\pi x}{L}\, dx \\
        b_n & =\frac{1}{L}\int_{-L}^{L}f(x)\sin\tfrac{n\pi x}{L}\, dx
    \end{align*}
\end{definition}

Intuitively, each coefficient \(a_n\) (resp.\ \(b_n\)) measures how much of the cosine (resp.\ sine) wave of frequency \(n\pi/L\) is present in \(f\).

The Fourier series can be understood geometrically as projecting the function \(f\) onto the orthogonal basis of trigonometric functions. Each coefficient represents the "component" of \(f\) in the direction of a particular basis function.

\begin{figure}[htbp]
    \centering
    \includestandalone[width=0.48\textwidth]{figures/fourier-ortho-1}
    \hfill
    \includestandalone[width=0.44\textwidth]{figures/fourier-ortho-2}
    \caption{
        Left: Decomposition of a function \(f(x)\) into its trigonometric components.
        Right: The orthogonal trigonometric basis functions on \([-L, L]\).
    }
    \label{fig:fourier-orthogonal-projection}
\end{figure}

The projection formula for coefficient \(a_n\) can be visualized as:
\begin{equation}
    a_n = \frac{\langle f, \cos(n\pi x/L) \rangle}{\|\cos(n\pi x/L)\|^2} = \frac{\int_{-L}^{L} f(x)\cos(n\pi x/L)\,dx}{\int_{-L}^{L} \cos^2(n\pi x/L)\,dx} = \frac{1}{L}\int_{-L}^{L} f(x)\cos(n\pi x/L)\,dx
\end{equation}

This geometric interpretation shows that Fourier analysis is fundamentally about finding the best approximation to \(f\) in the subspace spanned by trigonometric functions, where "best" means minimizing the \(L^2\) error over the interval \([-L,L]\).

\begin{definition}{Complex Fourier Series}{complex-fourier}
    The complex exponential form of the Fourier series is:
    \begin{equation}
        f(x) = \sum_{n=-\infty}^{\infty} c_n e^{i n \pi x / L}
    \end{equation}
    where the complex coefficients are:
    \begin{equation}
        c_n = \frac{1}{2L} \int_{-L}^{L} f(x) e^{-i n \pi x / L} \, dx
    \end{equation}

    The relationship between real and complex coefficients is:
    \begin{align}
        c_0    & = \frac{a_0}{2}                                        \\
        c_n    & = \frac{a_n - i b_n}{2} \quad (n > 0)                  \\
        c_{-n} & = \frac{a_n + i b_n}{2} = \overline{c_n} \quad (n > 0)
    \end{align}
\end{definition}

\subsection{Properties of Fourier Series}
The Fourier series possesses several important properties that mirror the structure of the underlying function space.
\begin{property}{Properties of Fourier Series}{fourier-properties}
    Let \( f(x) \) and \( g(x) \) be \( 2L \)-periodic functions with Fourier series, and let \( \alpha, \beta \) be constants. Then:
    \begin{enumerate}[label=\textbf{\arabic*.}, leftmargin=*, noitemsep]
        \item \textbf{Linearity:} The Fourier series is linear:
              \[
                  \mathcal{F}[\alpha f(x) + \beta g(x)] = \alpha \mathcal{F}[f(x)] + \beta \mathcal{F}[g(x)].
              \]
        \item \textbf{Differentiation:} If \( f \) is continuous and piecewise smooth, the Fourier series of \( f'(x) \) is obtained by term-wise differentiation:
              \[
                  f'(x) = \sum_{n=1}^{\infty} \left( -a_n \frac{n\pi}{L} \sin\left(\frac{n\pi x}{L}\right) + b_n \frac{n\pi}{L} \cos\left(\frac{n\pi x}{L}\right) \right),
              \]
              where \( a_n, b_n \) are the Fourier coefficients of \( f \).
        \item \textbf{Integration:} The Fourier series of \( f(x) \) can be integrated term by term:
              \[
                  \int_{-L}^{x} f(t)\,dt = \frac{a_0}{2}(x + L) + \sum_{n=1}^{\infty} \left[ a_n \frac{L}{n\pi} \sin\left(\frac{n\pi x}{L}\right) - b_n \frac{L}{n\pi} \cos\left(\frac{n\pi x}{L}\right) + b_n \frac{L}{n\pi} \right],
              \]
              where the integration constant is chosen so that the result is \( 2L \)-periodic.
    \end{enumerate}
\end{property}

\subsubsection{Parseval's Theorem}
Parseval's theorem relates the \(L^2\) norm of a function to the sum of the squares of its Fourier coefficients, providing a powerful tool for analyzing the energy content of signals.
\begin{theorem}{Parseval's Theorem}{parseval}
    Let \( f \) be a \( 2L \)-periodic, square-integrable function on \([-L, L]\), with complex Fourier coefficients
    \[
        c_n = \frac{1}{2L} \int_{-L}^{L} f(x) e^{-i n \pi x / L} \, dx.
    \]
    Define the inner product and norm on \( \mathrm{L}^2([-L, L]) \) by
    \[
        \langle f, g \rangle = \int_{-L}^{L} f(x) \cdot \overline{g(x)} \, dx, \qquad \|f\|^2 = \langle f, f \rangle = \int_{-L}^{L} |f(x)|^2 \, dx.
    \]
    Then Parseval's theorem states:
    \[
        \|f\|^2 = \int_{-L}^{L} |f(x)|^2 \, dx = 2L \sum_{n=-\infty}^{\infty} |c_n|^2.
    \]
    This expresses the completeness and orthogonality of the exponential basis in \( L^2([-L, L]) \).
\end{theorem}

\begin{proof}
    Let \( f(x) = \sum_{n=-\infty}^{\infty} c_n e^{i n \pi x / L} \) be the complex Fourier series of \( f \), converging in \( L^2([-L, L]) \).
    The coefficients are given by
    \[
        c_n = \frac{1}{2L} \int_{-L}^{L} f(x) e^{-i n \pi x / L} \, dx.
    \]
    The exponential functions \( \{e^{i n \pi x / L}\}_{n \in \mathbb{Z}} \) form an orthogonal basis for \( L^2([-L, L]) \) with
    \[
        \int_{-L}^{L} e^{i n \pi x / L} \overline{e^{i m \pi x / L}} \, dx = \int_{-L}^{L} e^{i (n-m) \pi x / L} \, dx =
        \begin{cases}
            2L, & n = m,    \\
            0,  & n \neq m.
        \end{cases}
    \]
    Therefore,
    \begin{align*}
        \|f\|^2 & = \int_{-L}^{L} |f(x)|^2 \, dx                                                                                                                                 \\
                & = \int_{-L}^{L} \left( \sum_{n=-\infty}^{\infty} c_n e^{i n \pi x / L} \right) \overline{ \left(\sum_{m=-\infty}^{\infty} c_m e^{i m \pi x / L} \right)} \, dx \\
                & = \int_{-L}^{L} \sum_{n} \sum_{m} c_n \overline{c_m} e^{i (n-m) \pi x / L} \, dx                                                                               \\
                & = \sum_{n} \sum_{m} c_n \overline{c_m} \int_{-L}^{L} e^{i (n-m) \pi x / L} \, dx                                                                               \\
                & = \sum_{n} \sum_{m} c_n \overline{c_m} \cdot 2L \, \delta_{nm}                                                                                                 \\
                & = 2L \sum_{n=-\infty}^{\infty} |c_n|^2.
    \end{align*}
    This completes the proof.
\end{proof}

\subsubsection{Gibbs Phenomenon}
When approximating a function with discontinuities using Fourier series, the partial sums exhibit an overshoot near the discontinuities that does not disappear as more terms are added. This overshoot is approximately 9\% of the jump discontinuity and is known as the \textbf{Gibbs phenomenon}.

\begin{figure}[H]
    \centering
    \includestandalone[width=0.8\textwidth]{figures/fourier-gibbs}
    \caption{Gibbs phenomenon showing persistent overshoot near discontinuities}
\end{figure}

\subsubsection{Convergence of Fourier Series}
A natural question is: \emph{When does the Fourier series of a function converge to the function itself?}
Not every periodic function has a Fourier series that converges everywhere to the original function.
The answer is provided by the \textbf{Dirichlet conditions}, which give sufficient (but not necessary) criteria for pointwise convergence of the Fourier series.

\begin{theorem}{Dirichlet Conditions for Fourier Series Convergence}{dirichlet-conditions}
    Let \( f \) be a \( 2L \)-periodic, piecewise smooth function on \(\mathbb{R}\). If the following conditions hold:
    \begin{enumerate}
        \item \( f \) is absolutely integrable over one period: \( \int_{-L}^{L} |f(x)|\,dx < \infty \).
        \item \( f \) has only finitely many discontinuities (all of finite size) and finitely many extrema in any period.
    \end{enumerate}
    Then, at every point \( x \) where \( f \) is continuous, its Fourier series converges to \( f(x) \). At a point of discontinuity \( x_0 \), the Fourier series converges to the average of the right and left limits:
    \begin{equation}
        \lim_{N \to \infty} S_N(x_0) = \frac{f(x_0^+) + f(x_0^-)}{2}
    \end{equation}
    where \( S_N(x) \) is the \( N \)-th partial sum of the Fourier series, and \( f(x_0^+), f(x_0^-) \) are the right and left limits of \( f \) at \( x_0 \).
\end{theorem}

\begin{itemize}
    \item Dirichlet conditions are sufficient (not necessary) for pointwise convergence.
    \item At discontinuities, the series converges to the average of left and right limits.
    \item Convergence is pointwise, not uniform; Gibbs phenomenon may appear near jumps.
\end{itemize}

\subsection{Examples of Fourier Series}
\begin{example}{Square Wave}{square-wave}
    Consider the square wave function defined on
    \begin{equation}
        f(x) = \begin{cases}
            1  & \text{if } 0 < x < \pi   \\
            -1 & \text{if } -\pi < x < 0  \\
            0  & \text{if } x = 0, \pm\pi
        \end{cases}
    \end{equation}

    Since the function is odd, all cosine terms vanish (\(a_n = 0\)), and we only have sine terms:
    \begin{align}
        b_n & = \frac{1}{\pi} \int_{-\pi}^{\pi} f(x) \sin(nx) \, dx    \\
            & = \frac{2}{\pi} \int_{0}^{\pi} \sin(nx) \, dx            \\
            & = \frac{2}{\pi} \left[-\frac{\cos(nx)}{n}\right]_0^{\pi} \\
            & = \frac{2}{n\pi}(1 - \cos(n\pi))                         \\
            & = \begin{cases}
                    \frac{4}{n\pi} & \text{if } n \text{ is odd}  \\
                    0              & \text{if } n \text{ is even}
                \end{cases}
    \end{align}

    Therefore, the Fourier series is:
    \begin{equation}
        f(x) = \frac{4}{\pi} \sum_{n=1,3,5,\ldots} \frac{\sin(nx)}{n} = \frac{4}{\pi}\left(\sin x + \frac{\sin 3x}{3} + \frac{\sin 5x}{5} + \cdots\right)
    \end{equation}
\end{example}

\begin{figure}[H]
    \centering
    \includestandalone[width=0.8\textwidth]{figures/fourier-square}
    \caption{Fourier series approximation of a square wave showing convergence with increasing number of terms}
\end{figure}

\begin{example}{Sawtooth Wave}{sawtooth-wave}
    Consider the sawtooth wave defined on \((-\pi, \pi)\):
    \begin{equation}
        f(x) = x \quad \text{for } x \in (-\pi, \pi)
    \end{equation}

    Since this function is odd, we have \(a_n = 0\) for all \(n\), and:
    \begin{align}
        b_n & = \frac{1}{\pi} \int_{-\pi}^{\pi} x \sin(nx) \, dx \\
            & = \frac{2}{\pi} \int_{0}^{\pi} x \sin(nx) \, dx    \\
            & = \frac{2(-1)^{n+1}}{n}
    \end{align}

    The Fourier series is:
    \begin{equation}
        f(x) = 2 \sum_{n=1}^{\infty} \frac{(-1)^{n+1}}{n} \sin(nx) = 2\left(\sin x - \frac{\sin 2x}{2} + \frac{\sin 3x}{3} - \cdots\right)
    \end{equation}
\end{example}

\begin{figure}[H]
    \centering
    \includestandalone[width=0.8\textwidth]{figures/fourier-sawtooth}
    \caption{Fourier series approximation of a sawtooth wave}
    \label{fig:fourier-sawtooth}
\end{figure}

\section{Discrete Fourier Transform (DFT)}

\subsection{Motivation: From Fourier Series to DFT}
The Discrete Fourier Transform (DFT) arises naturally when we sample a periodic function at finitely many points. While the classical Fourier series decomposes a function into sines and cosines (or complex exponentials) over a continuous interval, in practice we often work with discrete data: sampled signals, digital images, or numerical solutions. The DFT provides a way to analyze the frequency content of such discrete data, mirroring the role of the Fourier series for continuous functions.

\subsection{Definition and Properties of the DFT}
\begin{definition}{Discrete Fourier Transform}{dft}
    Given a sequence $\{f_k\}_{k=0}^{N-1}$ (e.g., $N$ samples of a periodic function), the DFT produces a sequence $\{F_n\}_{n=0}^{N-1}$:
    \begin{equation}
        F_n = \sum_{k=0}^{N-1} f_k \, e^{-2\pi i n k / N}, \qquad n = 0, 1, \ldots, N-1.
    \end{equation}
    The inverse DFT reconstructs the original sequence:
    \begin{equation}
        f_k = \frac{1}{N} \sum_{n=0}^{N-1} F_n \, e^{2\pi i n k / N}, \qquad k = 0, 1, \ldots, N-1.
    \end{equation}
\end{definition}

The DFT can be written as a matrix-vector multiplication:
\begin{equation}
    \mathbf{F} = \mathcal{F}_N \mathbf{f}
\end{equation}
where $\mathcal{F}_N$ is the $N \times N$ DFT matrix with entries $e^{-2\pi i jk / N}$.

\begin{definition}{DFT Matrix}{dft-matrix}
    The DFT matrix $\mathcal{F}_N$ is defined by
    \begin{equation}
        (\mathcal{F}_N)_{jk} = e^{-2\pi i jk / N} = \omega_N^{jk}, \qquad \omega_N = e^{-2\pi i / N}, \quad j, k = 0, \ldots, N-1.
    \end{equation}
    This matrix is a Vandermonde matrix built from the $N$th roots of unity.
\end{definition}

\paragraph{Key Properties}
\begin{itemize}
    \item \textbf{Linearity:} The DFT is a linear transformation.
    \item \textbf{Periodicity:} Both input and output sequences are periodic with period $N$.
    \item \textbf{Orthogonality:} The complex exponentials $e^{2\pi i n k / N}$ form an orthogonal basis for $\mathbb{C}^N$.
    \item \textbf{Parseval's Theorem (Discrete):}
          \[
              \sum_{k=0}^{N-1} |f_k|^2 = \frac{1}{N} \sum_{n=0}^{N-1} |F_n|^2
          \]
\end{itemize}

\subsection{Fast Fourier Transform (FFT)}

For large $N$, direct computation of the DFT via matrix multiplication requires $O(N^2)$ operations. The \textbf{Fast Fourier Transform (FFT)} is a family of algorithms that compute the DFT much more efficiently, in $O(N \log N)$ time, by exploiting symmetries in the roots of unity.

\subsubsection{Recursive Divide-and-Conquer Structure}
The most common FFT algorithm (Cooley--Tukey) recursively splits the DFT of a sequence of length $N$ (assumed even) into two DFTs of length $N/2$: one for the even-indexed elements and one for the odd-indexed elements. These are then combined using so-called \emph{butterfly operations}.

\begin{algorithm}[H]
    \caption{Radix-2 Cooley--Tukey FFT Algorithm}
    \begin{algorithmic}[1]
        \Require Sequence $f = (f_0, \ldots, f_{N-1})$ of length $N = 2^m$
        \Ensure DFT $F = (F_0, \ldots, F_{N-1})$
        \Function{FFT}{$f$}
        \If{$N = 1$}
        \State \Return $f$
        \Else
        \State Split $f$ into even and odd indexed parts:
        \State $f^{(\mathrm{even})}_k = f_{2k}$
        \State $f^{(\mathrm{odd})}_k = f_{2k+1}$
        \State $E \gets$ \Call{FFT}{$f^{(\text{even})}$}
        \State $O \gets$ \Call{FFT}{$f^{(\text{odd})}$}
        \For{$n = 0$ \textbf{to} $N/2-1$}
        \State $\omega \gets e^{-2\pi i n / N}$ \hfill (\textbf{twiddle factor})
        \State $F_n \gets E_n + \omega \cdot O_n$
        \State $F_{n+N/2} \gets E_n - \omega \cdot O_n$
        \EndFor
        \State \Return $F$
        \EndIf
        \EndFunction
    \end{algorithmic}
\end{algorithm}

\paragraph{Twiddle Factor} The \textbf{twiddle factor} $e^{-2\pi i n / N}$ is the complex exponential that introduces the necessary phase shift when combining the even and odd DFTs.

\paragraph{Butterfly Operation} The basic computation in the FFT is the \emph{butterfly operation}, which combines two values (from the even and odd DFTs) using the twiddle factor:
\[
    \begin{pmatrix}
        F_n \\
        F_{n+N/2}
    \end{pmatrix}
    =
    \begin{pmatrix}
        1 & \mathcal{F}_N^n  \\
        1 & -\mathcal{F}_N^n
    \end{pmatrix}
    \begin{pmatrix}
        E_n \\
        O_n
    \end{pmatrix}
\]
for $n = 0, \ldots, N/2-1$.

\paragraph{Complexity and Efficiency}
The FFT reduces the computational complexity of the DFT from $O(N^2)$ to $O(N \log N)$ by recursively breaking the problem into smaller subproblems and combining their results efficiently.

\subsection{Examples and Interpretation}

\begin{example}{DFT Matrix for $N=4$}{dft-matrix-4}
    For $N = 4$, the primitive 4th root of unity is $\omega_4 = e^{-2\pi i / 4} = i$. The DFT matrix is
    \[
        \mathcal{F}_4 =
        \begin{pmatrix}
            1 & 1  & 1  & 1  \\
            1 & i  & -1 & -i \\
            1 & -1 & 1  & -1 \\
            1 & -i & -1 & i
        \end{pmatrix}
    \]
    Thus, for $\mathbf{f} = (f_0, f_1, f_2, f_3)^\top$, the DFT is $\mathbf{F} = \mathcal{F}_4 \mathbf{f}$.
\end{example}

\begin{example}{DFT of a Simple Sequence}{dft-example}
    Let $N = 4$ and $f_0 = 1, f_1 = 0, f_2 = -1, f_3 = 0$. Compute the DFT:
    \begin{align*}
        F_0 & = 1 + 0 + (-1) + 0 = 0                                                                           \\
        F_1 & = 1 + 0 \cdot e^{-i\pi/2} + (-1) \cdot e^{-i\pi} + 0 \cdot e^{-i3\pi/2} = 1 + 0 + 1 + 0 = 2      \\
        F_2 & = 1 + 0 \cdot e^{-i\pi} + (-1) \cdot e^{-i2\pi} + 0 \cdot e^{-i3\pi} = 1 + 0 - 1 + 0 = 0         \\
        F_3 & = 1 + 0 \cdot e^{-i3\pi/2} + (-1) \cdot e^{-i9\pi/2} + 0 \cdot e^{-i27\pi/2} = 1 + 0 + 1 + 0 = 2
    \end{align*}
    Thus, the DFT is $ (0, 2, 0, 2) $.
\end{example}

\begin{example}{DFT Example: $s = (1, 3, 1, -1)$}{dft-example2}
    Let $N = 4$ and $s = (1, 3, 1, -1) \in \mathbb{R}^4$. The DFT is
    \[
        \hat{s}_k = \sum_{n=0}^{3} s_n e^{-2\pi i k n / 4}, \qquad k = 0,1,2,3.
    \]
    Compute each component:
    \begin{align*}
        \hat{s}_0 & = 1 + 3 + 1 + (-1) = 4                                  \\
        \hat{s}_1 & = 1 + 3 e^{-i\pi/2} + 1 e^{-i\pi} + (-1) e^{-i3\pi/2}   \\
                  & = 1 + 3i - 1 + i = 4i                                   \\
        \hat{s}_2 & = 1 + 3 e^{-i\pi} + 1 e^{-i2\pi} + (-1) e^{-i3\pi}      \\
                  & = 1 - 3 + 1 + 1 = 0                                     \\
        \hat{s}_3 & = 1 + 3 e^{-i3\pi/2} + 1 e^{-i3\pi} + (-1) e^{-i9\pi/2} \\
                  & = 1 - 3i - 1 - i = -4i
    \end{align*}
    So, the DFT is $\hat{s} = (4, 4i, 0, -4i)$.
\end{example}

\paragraph{Interpretation: Relation to Continuous Fourier Coefficients}
The DFT coefficients $\hat{s}_k$ are discrete analogues of the continuous Fourier complex coefficients $c_k(f)$. For a sampled $2L$-periodic function $f(x)$ at $N$ equispaced points, the DFT coefficients approximate $c_k(f)$ up to normalization and scaling. As $N$ increases and the sampling becomes finer, the DFT coefficients converge (after normalization) to the continuous Fourier coefficients.

\subsection{Applications}
DFT and FFT are fundamental tools in:
\begin{itemize}
    \item Signal and image processing (filtering, compression)
    \item Numerical solutions of differential equations (spectral methods)
    \item Data analysis and pattern recognition
\end{itemize}

\subsection{Further Examples and Discussion}

\subsubsection*{DFT Example: $s = (1, 3, 1, -1)$}

Let $N = 4$ and $s = (1, 3, 1, -1) \in \mathbb{R}^4$. The DFT is
\[
    \hat{s}_k = \sum_{n=0}^{3} s_n e^{-2\pi i k n / 4}, \qquad k = 0,1,2,3.
\]
Compute each component:
\begin{align*}
    \hat{s}_0 & = 1 + 3 + 1 + (-1) = 4                                  \\
    \hat{s}_1 & = 1 + 3 e^{-i\pi/2} + 1 e^{-i\pi} + (-1) e^{-i3\pi/2}   \\
              & = 1 + 3i - 1 + i = 4i                                   \\
    \hat{s}_2 & = 1 + 3 e^{-i\pi} + 1 e^{-i2\pi} + (-1) e^{-i3\pi}      \\
              & = 1 - 3 + 1 + 1 = 0                                     \\
    \hat{s}_3 & = 1 + 3 e^{-i3\pi/2} + 1 e^{-i3\pi} + (-1) e^{-i9\pi/2} \\
              & = 1 - 3i - 1 - i = -4i
\end{align*}
So, the DFT is $\hat{s} = (4, 4i, 0, -4i)$.

\paragraph{Relation to Fourier Complex Coefficients}
The DFT coefficients $\hat{s}_k$ are discrete analogues of the continuous Fourier complex coefficients $c_k(f)$. For a sampled $2L$-periodic function $f(x)$ at $N$ equispaced points, the DFT coefficients approximate $c_k(f)$ up to normalization and scaling. Specifically, as $N$ increases and the sampling becomes finer, the DFT coefficients converge (after normalization) to the continuous Fourier coefficients.

\subsubsection*{DFT as Matrix Multiplication and FFT Complexity}
The DFT can be written as a matrix-vector multiplication:
\[
    \hat{s} = F_N s
\]
where $F_N$ is the $N \times N$ matrix with entries $(F_N)_{jk} = e^{-2\pi i jk / N}$.
A general matrix-vector multiplication requires $O(N^2)$ operations. However, the DFT matrix has a special structure: its entries are powers of the primitive $N$th root of unity. This allows the DFT to be computed recursively by dividing the problem into smaller DFTs of size $N/2$ (for $N$ a power of $2$), combining the results using the so-called "butterfly" operations.
This recursive approach is the basis of the Fast Fourier Transform (FFT) algorithm, which reduces the computational complexity to $O(N \log N)$. The FFT exploits the symmetry and periodicity of the complex exponentials, making it much faster than a naive matrix multiplication for large $N$.